\documentclass{amsart}
% For dealing with references we use the comment environment
\usepackage{verbatim}
\newenvironment{reference}{\comment}{\endcomment}
%\newenvironment{reference}{}{}
\newenvironment{slogan}{\comment}{\endcomment}
\newenvironment{history}{\comment}{\endcomment}

% For commutative diagrams we use Xy-pic
\usepackage[all]{xy}

% We use 2cell for 2-commutative diagrams.
\xyoption{2cell}
\UseAllTwocells

% We use multicol for the list of chapters between chapters
\usepackage{multicol}

% This is generally recommended for better output
\usepackage{lmodern}
\usepackage[T1]{fontenc}

% For cross-file-references

% Package for hypertext links:
\usepackage{hyperref}

% For any local file, say "hello.tex" you want to link to please

% Theorem environments.
%
\theoremstyle{plain}
\newtheorem{theorem}[subsection]{Theorem}
\newtheorem{proposition}[subsection]{Proposition}
\newtheorem{lemma}[subsection]{Lemma}

\theoremstyle{definition}
\newtheorem{definition}[subsection]{Definition}
\newtheorem{example}[subsection]{Example}
\newtheorem{exercise}[subsection]{Exercise}
\newtheorem{situation}[subsection]{Situation}

\theoremstyle{remark}
\newtheorem{remark}[subsection]{Remark}
\newtheorem{remarks}[subsection]{Remarks}

\numberwithin{equation}{subsection}

% Macros
%
\def\lim{\mathop{\mathrm{lim}}\nolimits}
\def\colim{\mathop{\mathrm{colim}}\nolimits}
\def\Spec{\mathop{\mathrm{Spec}}}
\def\Hom{\mathop{\mathrm{Hom}}\nolimits}
\def\Ext{\mathop{\mathrm{Ext}}\nolimits}
\def\SheafHom{\mathop{\mathcal{H}\!\mathit{om}}\nolimits}
\def\SheafExt{\mathop{\mathcal{E}\!\mathit{xt}}\nolimits}
\def\Sch{\mathit{Sch}}
\def\Mor{\mathop{\mathrm{Mor}}\nolimits}
\def\Ob{\mathop{\mathrm{Ob}}\nolimits}
\def\Sh{\mathop{\mathit{Sh}}\nolimits}
\def\NL{\mathop{N\!L}\nolimits}
\def\CH{\mathop{\mathrm{CH}}\nolimits}
\def\proetale{{pro\text{-}\acute{e}tale}}
\def\etale{{\acute{e}tale}}
\def\QCoh{\mathit{QCoh}}
\def\Ker{\mathop{\mathrm{Ker}}}
\def\Im{\mathop{\mathrm{Im}}}
\def\Coker{\mathop{\mathrm{Coker}}}
\def\Coim{\mathop{\mathrm{Coim}}}

% Boxtimes
%
\DeclareMathSymbol{\boxtimes}{\mathbin}{AMSa}{"02}

%
% Macros for moduli stacks/spaces
%
\def\QCohstack{\mathcal{QC}\!\mathit{oh}}
\def\Cohstack{\mathcal{C}\!\mathit{oh}}
\def\Spacesstack{\mathcal{S}\!\mathit{paces}}
\def\Quotfunctor{\mathrm{Quot}}
\def\Hilbfunctor{\mathrm{Hilb}}
\def\Curvesstack{\mathcal{C}\!\mathit{urves}}
\def\Polarizedstack{\mathcal{P}\!\mathit{olarized}}
\def\Complexesstack{\mathcal{C}\!\mathit{omplexes}}
% \Pic is the operator that assigns to X its picard group, usage \Pic(X)
% \Picardstack_{X/B} denotes the Picard stack of X over B
% \Picardfunctor_{X/B} denotes the Picard functor of X over B
\def\Pic{\mathop{\mathrm{Pic}}\nolimits}
\def\Picardstack{\mathcal{P}\!\mathit{ic}}
\def\Picardfunctor{\mathrm{Pic}}
\def\Deformationcategory{\mathcal{D}\!\mathit{ef}}

\setlength{\emergencystretch}{3em}
\begin{document}
\title[Stacks Project, Tag 05R0]{359 --- The nine-triangle diagram in a triangulated category}
\maketitle
\noindent Copyright (C) 2005--2025 Johan de Jong. Stacks Project authors. Source: \href{https://stacks.math.columbia.edu/tag/05R0}{Tag 05R0}.

\noindent Editorial difficulty note (not author text). Difficulty H2: The proof applies the octahedral axiom three times to construct all rows and columns. It then checks the factorization identities controlling the central squares and the required anticommutativity at the lower-right corner. This is a substantial sign-sensitive existence theorem from the abstract triangulated axioms..

\noindent Editorial provenance note (not author text). History: excerpt from revision \texttt{a0444\allowbreak 6e57e\allowbreak c1fbc\allowbreak 252a8\allowbreak 71afc\allowbreak ec775\allowbreak 2fb28\allowbreak 07b14}, derived.tex, lines 1030--1140. Reference presentation was adapted; original-author context and core support blocks were retained or added. The mathematical wording is unchanged. Licensed under GNU FDL 1.2 or later; no invariant sections or cover texts. See ../licenses/COPYING. Context blocks reproduce author definitions and supporting results; they are not additional counted problems.

\begin{definition}
\label{definition-triangulated-category}
A {\it triangulated category} consists of a triple
$(\mathcal{D}, \{[n]\}_{n\in \mathbf{Z}}, \mathcal{T})$
where
\begin{enumerate}
\item $\mathcal{D}$ is an additive category,
\item $[1] : \mathcal{D} \to \mathcal{D}$, $E \mapsto E[1]$
is an additive auto-equivalence and $[n]$ for $n \in \mathbf{Z}$ is
as discussed above, and
\item $\mathcal{T}$ is a set of triangles (Definition \ref{definition-triangle})
called the {\it distinguished triangles}
\end{enumerate}
subject to the following conditions
\begin{enumerate}
\item[TR1] Any triangle isomorphic to a distinguished triangle is
a distinguished triangle. Any triangle of the form
$(X, X, 0, \text{id}, 0, 0)$ is distinguished.
For any morphism $f : X \to Y$ of $\mathcal{D}$ there exists a
distinguished triangle of the form $(X, Y, Z, f, g, h)$.
\item[TR2] The triangle $(X, Y, Z, f, g, h)$ is distinguished
if and only if the triangle $(Y, Z, X[1], g, h, -f[1])$ is.
\item[TR3] Given a solid diagram
$$
\xymatrix{
X \ar[r]^f \ar[d]^a &
Y \ar[r]^g \ar[d]^b &
Z \ar[r]^h \ar@{-->}[d] &
X[1] \ar[d]^{a[1]} \\
X' \ar[r]^{f'} &
Y' \ar[r]^{g'} &
Z' \ar[r]^{h'} &
X'[1]
}
$$
whose rows are distinguished triangles and which satisfies
$b \circ f = f' \circ a$, there exists a morphism
$c : Z \to Z'$ such that $(a, b, c)$ is a morphism of triangles.
\item[TR4] Given objects $X$, $Y$, $Z$ of $\mathcal{D}$, and morphisms
$f : X \to Y$, $g : Y \to Z$, and distinguished triangles
$(X, Y, Q_1, f, p_1, d_1)$,
$(X, Z, Q_2, g \circ f, p_2, d_2)$,
and
$(Y, Z, Q_3, g, p_3, d_3)$,
there exist
morphisms $a : Q_1 \to Q_2$ and $b : Q_2 \to Q_3$ such
that
\begin{enumerate}
\item $(Q_1, Q_2, Q_3, a, b, p_1[1] \circ d_3)$ is a
distinguished triangle,
\item the triple $(\text{id}_X, g, a)$ is
a morphism of triangles
$(X, Y, Q_1, f, p_1, d_1) \to (X, Z, Q_2, g \circ f, p_2, d_2)$, and
\item the triple $(f, \text{id}_Z, b)$ is a morphism of triangles
$(X, Z, Q_2, g \circ f, p_2, d_2) \to (Y, Z, Q_3, g, p_3, d_3)$.
\end{enumerate}
\end{enumerate}
We will call $(\mathcal{D}, [\ ], \mathcal{T})$ a
{\it pre-triangulated category} if TR1, TR2 and TR3
hold.\footnote{We use $[\ ]$ as an abbreviation for the
family $\{[n]\}_{n\in \mathbf{Z}}$.}
\end{definition}

\begin{definition}
\label{definition-triangle}
Let $\mathcal{D}$ be an additive category. Let
$[1] : \mathcal{D} \to \mathcal{D}$, $E \mapsto E[1]$
be an additive functor which is an auto-equivalence of $\mathcal{D}$.
\begin{enumerate}
\item A {\it triangle} is a sextuple
$(X, Y, Z, f, g, h)$ where $X, Y, Z \in \Ob(\mathcal{D})$ and
$f : X \to Y$, $g : Y \to Z$ and $h : Z \to X[1]$ are morphisms
of $\mathcal{D}$.
\item A {\it morphism of triangles}
$(X, Y, Z, f, g, h) \to (X', Y', Z', f', g', h')$
is given by morphisms $a : X \to X'$, $b : Y \to Y'$ and $c : Z \to Z'$
of $\mathcal{D}$ such that
$b \circ f = f' \circ a$, $c  \circ g = g' \circ b$ and
$a[1] \circ h = h' \circ c$.
\end{enumerate}
\end{definition}

\begin{proposition}
\label{proposition-9}
Let $\mathcal{D}$ be a triangulated category. Any commutative diagram
$$
\xymatrix{
X \ar[r] \ar[d] & Y \ar[d] \\
X' \ar[r] & Y'
}
$$
can be extended to a diagram
$$
\xymatrix{
X \ar[r] \ar[d] & Y \ar[r] \ar[d] & Z \ar[r] \ar[d] & X[1] \ar[d] \\
X' \ar[r] \ar[d] & Y' \ar[r] \ar[d] & Z' \ar[r] \ar[d] & X'[1] \ar[d] \\
X'' \ar[r] \ar[d] & Y'' \ar[r] \ar[d] & Z'' \ar[r] \ar[d] & X''[1] \ar[d] \\
X[1] \ar[r] & Y[1] \ar[r] & Z[1] \ar[r] & X[2]
}
$$
where all the squares are commutative, except for the lower right square
which is anticommutative. Moreover, each of the rows and columns are
distinguished triangles. Finally, the morphisms on the bottom row
(resp.\ right column) are obtained from the morphisms of the top row
(resp.\ left column) by applying $[1]$.
\end{proposition}

\begin{proof}
During this proof we avoid writing the arrows in order to make the proof
legible. Choose distinguished triangles
$(X, Y, Z)$, $(X', Y', Z')$, $(X, X', X'')$, $(Y, Y', Y'')$, and
$(X, Y', A)$. Note that the morphism $X \to Y'$ is both equal
to the composition $X \to Y \to Y'$ and equal to the composition
$X \to X' \to Y'$. Hence, we can find morphisms
\begin{enumerate}
\item $a : Z \to A$ and $b : A \to Y''$, and
\item $a' : X'' \to A$ and $b' : A \to Z'$
\end{enumerate}
as in TR4. Denote $c : Y'' \to Z[1]$ the composition
$Y'' \to Y[1] \to Z[1]$ and denote $c' : Z' \to X''[1]$ the composition
$Z' \to X'[1] \to X''[1]$. The conclusion of our application TR4
are that
\begin{enumerate}
\item $(Z, A, Y'', a, b, c)$, $(X'', A, Z', a', b', c')$
are distinguished triangles,
\item $(X, Y, Z) \to (X, Y', A)$,
$(X, Y', A) \to (Y, Y', Y'')$,
$(X, X', X'') \to (X, Y', A)$,
$(X, Y', A) \to (X', Y', Z')$
are morphisms of triangles.
\end{enumerate}
First using that
$(X, X', X'') \to (X, Y', A)$ and $(X, Y', A) \to (Y, Y', Y'')$.
are morphisms of triangles we see the first of the diagrams
$$
\vcenter{
\xymatrix{
X' \ar[r] \ar[d] & Y' \ar[d] \\
X'' \ar[r]^{b \circ a'} \ar[d] & Y'' \ar[d] \\
X[1] \ar[r] & Y[1]
}
}
\quad\text{and}\quad
\vcenter{
\xymatrix{
Y \ar[r] \ar[d] & Z \ar[d]^{b' \circ a} \ar[r] & X[1] \ar[d] \\
Y' \ar[r] & Z' \ar[r] & X'[1]
}
}
$$
is commutative. The second is commutative too using that
$(X, Y, Z) \to (X, Y', A)$ and $(X, Y', A) \to (X', Y', Z')$ are morphisms
of triangles. At this point we choose a distinguished triangle
$(X'', Y'' , Z'')$ starting with the map $b \circ a' : X'' \to Y''$.

\medskip\noindent
Next we apply TR4 one more time to the morphisms
$X'' \to A \to Y''$ and the triangles
$(X'', A, Z', a', b', c')$,
$(X'', Y'', Z'')$, and
$(A, Y'', Z[1], b, c , -a[1])$ to get morphisms
$a'' : Z' \to Z''$ and $b'' : Z'' \to Z[1]$.
Then $(Z', Z'', Z[1], a'', b'', - b'[1] \circ a[1])$ is a distinguished
triangle, hence also $(Z, Z', Z'', -b' \circ a, a'', -b'')$
and hence also $(Z, Z', Z'', b' \circ a, a'', b'')$.
Moreover, $(X'', A, Z') \to (X'', Y'', Z'')$ and
$(X'', Y'', Z'') \to (A, Y'', Z[1], b, c , -a[1])$
are morphisms of triangles.
At this point we have defined all the distinguished triangles
and all the morphisms, and all that's left is to verify some
commutativity relations.

\medskip\noindent
To see that the middle square in the diagram commutes, note
that the arrow $Y' \to Z'$ factors as $Y' \to A \to Z'$
because $(X, Y', A) \to (X', Y', Z')$ is a morphism of triangles.
Similarly, the morphism $Y' \to Y''$ factors as
$Y' \to A \to Y''$ because $(X, Y', A) \to (Y, Y', Y'')$ is a
morphism of triangles. Hence the middle square commutes because
the square with sides $(A, Z', Z'', Y'')$ commutes as
$(X'', A, Z') \to (X'', Y'', Z'')$ is a morphism of triangles (by TR4).
The square with sides $(Y'', Z'', Y[1], Z[1])$ commutes
because $(X'', Y'', Z'') \to (A, Y'', Z[1], b, c , -a[1])$
is a morphism of triangles and $c : Y'' \to Z[1]$ is the composition
$Y'' \to Y[1] \to Z[1]$.
The square with sides $(Z', X'[1], X''[1], Z'')$ is commutative
because $(X'', A, Z') \to (X'', Y'', Z'')$ is a morphism of triangles
and $c' : Z' \to X''[1]$ is the composition $Z' \to X'[1] \to X''[1]$.
Finally, we have to show that the square with sides
$(Z'', X''[1], Z[1], X[2])$ anticommutes. This holds because
$(X'', Y'', Z'') \to (A, Y'', Z[1], b, c , -a[1])$
is a morphism of triangles and we're done.
\end{proof}
\end{document}
